\documentclass[../../main.tex]{subfiles}
\begin{document}

\begin{center}
  {\fontsize{48pt}{58pt}\selectfont \textbf{1982}\normalsize\textbf{年度}}
\end{center}

\begin{table}[htb]
  \centering
  \caption{出題テーマと難易度}
  \begin{tabular}{|c|p{9cm}|c|}
    \hline
    \textbf{問題番号} & \textbf{テーマ} & \textbf{難易度} \\
    \hline
    第1問 & 三次曲線の接線と囲まれる面積比 & \\
    \hline
    第2問 & 二次方程式の根で定まる三角形領域 & \\
    \hline
    第3問 & 重心を順に結ぶ四辺形の収束 & \\
    \hline
    第4問 & 立方体の平面切断面積 & \\
    \hline
    第5問 & 硬貨の表裏による歩行と到達確率 & \\
    \hline
    第6問 & 指数関数を含む移動積分の極値 & \\
    \hline
  \end{tabular}
\end{table}

\begin{table}[htb]
  \centering
  \caption{解答}
  \renewcommand{\arraystretch}{1.4}
  \begin{tabular}{|c|c|p{9cm}|}
    \hline
    \textbf{問題番号} & \textbf{小問} & \textbf{答え} \\
    \hline
    第1問 & (1) & 接線は別の点で曲線と交わる（証明問題） \\
    \cline{2-3}
    & (2) & $P$の$x$座標が正なら面積比$8:1$，負なら$1:8$ \\
    \hline
    第2問 & -- & $\dfrac{2}{3}<a\le1$または$2\le a<3$ \\
    \hline
    第3問 & (1) & 4線分は1点$P$を共有する \\
    \cline{2-3}
    & (2) & $A_n$は一直線上にある \\
    \cline{2-3}
    & (3) & $\displaystyle\lim_{n\to\infty}A_nP=0$ \\
    \hline
    第4問 & (1) & $T(t)=\dfrac{\sqrt{3}}{2}t^2$ \\
    \cline{2-3}
    & (2) & 最大切断面積は$\dfrac{3\sqrt{3}}{4}$ \\
    \hline
    第5問 & (1) & $p_n=\dfrac{{}_{n+4}C_4}{2^{n+5}}$ \\
    \cline{2-3}
    & (2) & $n=3,4$ \\
    \hline
    第6問 & (1) & $-1\le x\le0$で$f(-1)\le f(x)\le f(0)$ \\
    \cline{2-3}
    & (2) & 最大は$x=\dfrac{1}{e-1}$，最小は$x=-\dfrac{e}{e-1}$ \\
    \hline
  \end{tabular}
\end{table}

\paragraph{試験形態}
6問．試験時間は要確認．

\end{document}
